% Výsledná formule převodu Sudoku na SAT. Každá část podmínky je zvýrazněna
% vlastní barvou a doplněna bublinou s vysvětlivkou.
\begin{tikzpicture}[
    font=\small,
    cast/.style={
        anchor=west,
        rounded corners=3pt,
        inner xsep=5pt,
        inner ysep=4pt,
        minimum height=11mm,
        line width=0.5pt
    },
    spojka/.style={anchor=east},
    bublina/.style={
        rectangle callout,
        anchor=west,
        rounded corners=3pt,
        line width=0.6pt,
        inner sep=4pt,
        text width=27mm,
        align=left,
        font=\scriptsize,
        xshift=7mm
    }
]

    % ------------------------------------------------------------------
    % Jednotlivé části formule
    % ------------------------------------------------------------------
    \node[cast, fill=green!8, draw=exampleblocktitlefg!60] (bunkyA) at (0.2,0)
        {$\displaystyle\bigwedge_{r,c}\;\bigvee_{d}x_{r,c,d}$};

    \node[cast, fill=red!8, draw=alertblocktitlefg!60] (bunkyN) at (0.2,-1.35)
        {$\displaystyle\bigwedge_{r,c}\;\bigwedge_{d\neq d^{\prime}}(\neg x_{r,c,d}\lor\neg x_{r,c,d^{\prime}})$};

    \node[cast, fill=blue!7, draw=darkblue!60] (radky) at (0.2,-2.9)
        {$\displaystyle\bigwedge_{r,d}\left[\bigvee_{c}x_{r,c,d}\;\land\bigwedge_{c\neq c^{\prime}}(\neg x_{r,c,d}\lor\neg x_{r,c^{\prime},d})\right]$};

    \node[cast, fill=yellow!12, draw=timportantboxcolor!60] (sloupce) at (0.2,-4.5)
        {$\displaystyle\bigwedge_{c,d}\left[\bigvee_{r}x_{r,c,d}\;\land\bigwedge_{r\neq r^{\prime}}(\neg x_{r,c,d}\lor\neg x_{r^{\prime},c,d})\right]$};

    \node[cast, fill=violet!8, draw=violet!60] (bloky) at (0.2,-6.15)
        {$\displaystyle\bigwedge_{B,d}\left[\bigvee_{(r,c)\in B}x_{r,c,d}\;\land\bigwedge_{\substack{(r,c),(r^{\prime},c^{\prime})\in B\\ (r,c)\neq(r^{\prime},c^{\prime})}}(\neg x_{r,c,d}\lor\neg x_{r^{\prime},c^{\prime},d})\right]$};

    \node[cast, fill=maincolor!8, draw=maincolor!60] (zadani) at (0.2,-7.75)
        {$\displaystyle\bigwedge_{(r,c,d)\in\mathcal{Z}}x_{r,c,d}$};

    % ------------------------------------------------------------------
    % Levý sloupec se spojkami
    % ------------------------------------------------------------------
    \node[spojka] at (0.1,0)     {$\varphi\;=$};
    \node[spojka] at (0.1,-1.35) {$\land$};
    \node[spojka] at (0.1,-2.9)  {$\land$};
    \node[spojka] at (0.1,-4.5) {$\land$};
    \node[spojka] at (0.1,-6.15)  {$\land$};
    \node[spojka] at (0.1,-7.75) {$\land$};

    % ------------------------------------------------------------------
    % Bubliny s vysvětlivkami (zarovnané k nejširší části formule)
    % ------------------------------------------------------------------
    \node[bublina, fill=green!8, draw=exampleblocktitlefg,
        callout absolute pointer={(bunkyA.east)}]
        at (bloky.east |- bunkyA) {Každá buňka obsahuje \textbf{alespoň} jednu číslici};

    \node[bublina, fill=red!8, draw=alertblocktitlefg,
        callout absolute pointer={(bunkyN.east)}]
        at (bloky.east |- bunkyN) {Každá buňka obsahuje \textbf{nejvýše} jednu číslici};

    \node[bublina, fill=blue!7, draw=darkblue,
        callout absolute pointer={(radky.east)}]
        at (bloky.east |- radky) {V~každém \textbf{řádku} je každá číslice právě jednou};

    \node[bublina, fill=yellow!12, draw=timportantboxcolor,
        callout absolute pointer={(sloupce.east)}]
        at (bloky.east |- sloupce) {V~každém \textbf{sloupci} je každá číslice právě jednou};

    \node[bublina, fill=violet!8, draw=violet,
        callout absolute pointer={(bloky.east)}]
        at (bloky.east |- bloky) {V~každém \textbf{bloku} je každá číslice právě jednou};

    \node[bublina, fill=maincolor!8, draw=maincolor,
        callout absolute pointer={(zadani.east)}]
        at (bloky.east |- zadani) {Buňky \textbf{zadané} v~počáteční tabulce};

\end{tikzpicture}
